
\subsection*{מטרות השיעור}
\begin{itemize}
    \item חזרה על סימונים ומושגי יסוד ממכניקה קוונטית 1
    \item הבנת אופרטורים נורמליים ותכונותיהם
    \item למידת ייצוג סימטריות באמצעות אופרטורים אוניטריים ואנטי-אוניטריים
    \item היכרות עם משפט ויגנר והשלכותיו הפיזיקליות
\end{itemize}

\subsection*{תזכורת: סימונים בסיסיים}
נזכיר תחילה את הסימונים המרכזיים שנשתמש בהם לאורך הקורס:
\begin{itemize}
    \item \textbf{מצבים קוונטיים:} $\ket{\psi}, \ket{\phi}$ — וקטורי מצב במרחב הילברט
    \item \textbf{מכפלה פנימית:} $\braket{\phi|\psi}$ — מכפלה סקלרית של מצבים
    \item \textbf{אופרטורים:} $\hat{A}, \hat{H}, \hat{p}$ — טרנספורמציות לינאריות על מרחב המצבים
    \item \textbf{אופרטור צמוד:} $\hat{A}^\dagger$ — מקיים $\bra{\phi}\hat{A}\ket{\psi} = \braket{\hat{A}^\dagger\phi|\psi}$
    \item \textbf{קומוטטור:} $[\hat{A},\hat{B}] = \hat{A}\hat{B} - \hat{B}\hat{A}$
\end{itemize}

\subsection{אופרטורים נורמליים}

\subsubsection*{הגדרה מרכזית}

\begin{definitionbox}{אופרטור נורמלי}
אופרטור $\hat{X}$ נקרא \textbf{נורמלי} אם הוא מתחלף עם הצמוד שלו:
\[
[\hat{X},\hat{X}^{\dagger}]=0 \quad \iff \quad \hat{X}^{\dagger}\hat{X} = \hat{X}\hat{X}^{\dagger}
\]
\end{definitionbox}

\subsubsection*{דוגמאות לאופרטורים נורמליים}

\begin{examplebox}{מקרים פרטיים חשובים}
\begin{enumerate}
    \item \textbf{אופרטור הזהות} $\hat{I}$ ו\textbf{אופרטור האפס} $\hat{0}$ — נורמליים באופן טריוויאלי
    
    \item \textbf{אופרטור הרמיטי} (צמוד עצמי): $\hat{H} = \hat{H}^\dagger$
    \begin{align*}
    [\hat{H},\hat{H}^\dagger] = [\hat{H},\hat{H}] = \hat{H}\hat{H} - \hat{H}\hat{H} = 0
    \end{align*}
    דוגמאות: המילטוניאן $\hat{H}$, אופרטור המיקום $\hat{x}$, אופרטור התנע $\hat{p}$
    
    \item \textbf{אופרטור אוניטרי}: $\hat{U}^{-1} = \hat{U}^\dagger$ (צמוד שווה להופכי)
    \begin{align*}
    [\hat{U},\hat{U}^\dagger] &= \hat{U}\hat{U}^\dagger - \hat{U}^\dagger\hat{U} \\
    &= \hat{U}\hat{U}^{-1} - \hat{U}^{-1}\hat{U} = \hat{I} - \hat{I} = 0
    \end{align*}
    דוגמאות: אופרטורי סיבוב, אופרטורי הזזה במרחב
\end{enumerate}
\end{examplebox}

\subsubsection*{תכונה מרכזית: לכסינות}

\begin{theorembox}{משפט הלכסון הספקטרלי}
כל אופרטור נורמלי $\hat{X}$ ניתן ללכסון על ידי בסיס אורתונורמלי של וקטורים עצמיים.

כלומר, קיים בסיס $\{\ket{\varphi_n}\}$ כך שלכל $n$:
\[
\hat{X}\ket{\varphi_n}=\lambda_n\ket{\varphi_n}
\]
כאשר $\lambda_n$ הם הערכים העצמיים (יכולים להיות מרוכבים), והוקטורים העצמיים מקיימים:
\[
\braket{\varphi_m|\varphi_n} = \delta_{mn}
\]
\end{theorembox}

\begin{notebox}{הערה חשובה}
עבור אופרטורים הרמיטיים, הערכים העצמיים $\lambda_n$ חייבים להיות \textbf{ממשיים}, והווקטורים העצמיים השייכים לערכים עצמיים שונים הם \textbf{אורתוגונליים}. זו תוצאה ישירה מכך שאופרטורים הרמיטיים הם מקרה פרטי של אופרטורים נורמליים.
\end{notebox}

\subsection{סימטריות במכניקה הקוונטית}

\subsubsection*{מושג הסימטריה}

\begin{definitionbox}{סימטריה קוונטית}
נאמר שמערכת קוונטית היא \textbf{סימטרית} תחת טרנספורמציה $T$ אם לכל שני מצבים $\ket{\psi}$ ו-$\ket{\phi}$ מתקיים:
\[
\left|\braket{\phi|\psi}\right|^{2}=\left|\braket{\phi'|\psi'}\right|^2
\]
כאשר $\ket{\psi'} = T\ket{\psi}$ ו-$\ket{\phi'} = T\ket{\phi}$ הם המצבים המותמרים.
\end{definitionbox}

\begin{notebox}{פרשנות פיזיקלית}
תנאי זה אומר ש\textbf{הסתברויות המעבר} בין מצבים נשמרות תחת הטרנספורמציה. 

אם $P_\phi = \ket{\phi}\bra{\phi}$ הוא אופרטור ההיטל על המצב $\ket{\phi}$, אז ההסתברות למצוא את המערכת במצב $\ket{\phi}$ כאשר היא נמצאת במצב $\ket{\psi}$ היא:
\[
P(\phi|\psi) = |\braket{\phi|\psi}|^2
\]
סימטריה דורשת ש-$P(\phi'|\psi') = P(\phi|\psi)$ — כלומר, הפיזיקה לא משתנה תחת הטרנספורמציה.
\end{notebox}

\subsubsection*{סימטריה חלקית \LR{vs.} סימטריה מלאה}

\begin{itemize}
    \item \textbf{סימטריה מלאה:} כל מצבי המערכת והמדידות נשארים בעלי אותה פיזיקה תחת הטרנספורמציה
    \item \textbf{סימטריה חלקית:} חלק מהמצבים משמרים הסתברויות ואחרים לא — במקרה זה לא נאמר שהמערכת כולה סימטרית
\end{itemize}

\subsubsection*{ייצוג סימטריות באמצעות אופרטורים}

כל סימטריה פיזיקלית $T$ מיוצגת על ידי אופרטור $\hat{U}_T$ הפועל על מרחב המצבים:
\[
\ket{\psi'} = \hat{U}_T\ket{\psi}
\]

\begin{examplebox}{דוגמאות לסימטריות}
\begin{enumerate}
    \item \textbf{סיבוב במרחב:} אם למערכת יש סימטריה ספרית (כמו אטום המימן), סיבוב הקואורדינטות לא משנה את הפיזיקה
    \[
    \ket{\psi'} = \hat{R}(\theta, \mathbf{n})\ket{\psi}
    \]
    כאשר $\hat{R}(\theta, \mathbf{n})$ הוא אופרטור סיבוב בזווית $\theta$ סביב ציר $\mathbf{n}$
    
    \item \textbf{הזזה במרחב:} אם המערכת הומוגנית, הזזה לא משנה את הפיזיקה
    \[
    \ket{\psi'} = \hat{T}(\mathbf{a})\ket{\psi}
    \]
    
    \item \textbf{שיקוף:} שיקוף מרחבי (\LR{parity}) $\hat{P}$ — מחליף $\mathbf{r} \to -\mathbf{r}$
\end{enumerate}
\end{examplebox}

\subsubsection*{קשר בין סימטריות לאופרטורים אוניטריים}

\begin{theorembox}{אופרטורי סימטריה}
כמעט כל סימטריה פיזיקלית מיוצגת על ידי אופרטור \textbf{אוניטרי} $\hat{U}$ המקיים:
\[
\hat{U}^\dagger \hat{U} = \hat{U} \hat{U}^\dagger = \hat{I}
\]

תכונה זו מבטיחה שמירת נורמות:
\[
\braket{\psi'|\psi'} = \bra{\psi}\hat{U}^\dagger\hat{U}\ket{\psi} = \braket{\psi|\psi}
\]
\end{theorembox}

\begin{notebox}{חריגה חשובה}
קיימת דוגמה חשובה אחת לסימטריה שאינה אוניטרית: \textbf{היפוך זמן} (\LR{Time Reversal}). 

כפי שנראה בהמשך, סימטריה זו מיוצגת על ידי אופרטור \textbf{אנטי-אוניטרי}.
\end{notebox}

\subsection{אופרטורים אנטי-לינאריים ואנטי-אוניטריים}

לפני שנציג את משפט ויגנר, עלינו להכיר סוג חדש של אופרטורים.

\subsubsection*{הגדרות}

\begin{definitionbox}{אופרטור אנטי-לינארי}
אופרטור $\hat{A}$ נקרא \textbf{אנטי-לינארי} אם הוא מקיים:
\[
\hat{A}(c_1\ket{\psi_1} + c_2\ket{\psi_2}) = c_1^*\hat{A}\ket{\psi_1} + c_2^*\hat{A}\ket{\psi_2}
\]
כלומר, הוא לינארי ביחס למצבים, אבל \textbf{מצמיד} את הסקלרים המרוכבים.
\end{definitionbox}

\begin{examplebox}{דוגמה: אופרטור הצימוד המרוכב}
נגדיר אופרטור $\hat{K}$ על ידי פעולתו על בסיס אורתונורמלי $\{\ket{n}\}$:
\[
\hat{K}\left(\sum_{n} c_n\ket{n}\right) = \sum_{n} c_n^*\ket{n}
\]

אופרטור זה הוא אנטי-לינארי. נבדוק:
\begin{align*}
\hat{K}(c_1\ket{\psi_1} + c_2\ket{\psi_2}) &= \hat{K}\left(\sum_n (c_1 a_n^{(1)} + c_2 a_n^{(2)})\ket{n}\right) \\
&= \sum_n (c_1 a_n^{(1)} + c_2 a_n^{(2)})^*\ket{n} \\
&= \sum_n (c_1^* a_n^{(1)*} + c_2^* a_n^{(2)*})\ket{n} \\
&= c_1^* \hat{K}\ket{\psi_1} + c_2^* \hat{K}\ket{\psi_2}
\end{align*}
\end{examplebox}

\begin{definitionbox}{אופרטור אוניטרי (תזכורת)}
אופרטור $\hat{U}$ הוא \textbf{אוניטרי} אם הוא ליניארי, הפיך, ומקיים:
\[
\braket{\hat{U}\psi|\hat{U}\phi} = \braket{\psi|\phi}
\]
כלומר, הוא שומר על המכפלה הפנימית (ועל הנורמה).
\end{definitionbox}

\begin{definitionbox}{אופרטור אנטי-אוניטרי}
אופרטור $\hat{A}$ הוא \textbf{אנטי-אוניטרי} אם הוא אנטי-לינארי, הפיך, ומקיים:
\[
\braket{\hat{A}\psi|\hat{A}\phi} = \braket{\phi|\psi}^*
\]
שימו לב להיפוך הסדר והצימוד המרוכב בצד ימין!
\end{definitionbox}

\subsubsection*{תכונות של אופרטורים אנטי-אוניטריים}

\begin{itemize}
    \item \textbf{שימור נורמות:} 
    \[
    \|\hat{A}\ket{\psi}\|^2 = \braket{\hat{A}\psi|\hat{A}\psi} = \braket{\psi|\psi}^* = \braket{\psi|\psi} = \|\ket{\psi}\|^2
    \]
    
    \item \textbf{הרכבה של שני אופרטורים אנטי-אוניטריים:}
    \[
    \hat{A}_1\hat{A}_2 \text{ הוא אופרטור לינארי ואוניטרי}
    \]
    הסיבה: הצימוד המרוכב מופעל פעמיים ומבטל את עצמו.
    
    \item \textbf{פירוק לאוניטרי וצימוד:} כל אופרטור אנטי-אוניטרי $\hat{A}$ ניתן לכתוב בצורה:
    \[
    \hat{A} = \hat{U}\hat{K}
    \]
    כאשר $\hat{U}$ אוניטרי ו-$\hat{K}$ הוא אופרטור הצימוד המרוכב.
\end{itemize}

\begin{notebox}{למה אנחנו צריכים את זה?}
אופרטורים אנטי-אוניטריים אינם מופיעים לעיתים קרובות בפיזיקה, אבל הם \textbf{הכרחיים} לתיאור \textbf{היפוך זמן}. 

הסיבה: היפוך זמן הופך $t \to -t$, ולכן צריך להפוך גם את $i \to -i$ במשוואת שרדינגר, מה שדורש אופרטור אנטי-לינארי.
\end{notebox}

\subsection*{עקרון ויגנר (\LR{Wigner's theorem})}

כל סימטריה במכניקה קוונטית — כלומר טרנספורמציה $T$ 

המשמרת הסתברויות מעבר בין מצבים — ניתנת למימוש על־ידי אופרטור אוניטרי
או אנטי־אוניטרי במרחב הילברט.

אם $T$ פועל על קרני המרחב כך ש-
\[
|\langle\psi|\phi\rangle|=|\langle T\psi|T\phi\rangle|,
\]
 אז קיימת טרנספורמציה $U_{T}$ שהיא אוניטרית או אנטי-אוניטרית, כך
ש-
\[
T|\psi\rangle=e^{i\alpha(\psi)}U_{T}|\psi\rangle
\]

\section*{דוגמות לסימטריות פיזיקליות}

\begin{center}
\begin{tabular}{|c|c|p{8cm}|}
\hline
\textbf{סימטריה} & \textbf{סוג האופרטור} & \textbf{תיאור} \\
\hline
סיבוב, הזזה, שיקוף במרחב & אוניטרי &
שומרים את מבנה המכפלה הפנימית באופן ליניארי. \\
\hline
היפוך זמן $T$ & אנטי־אוניטרי &
הופך את $i \mapsto -i$, ולכן חייב להיות אנטי־ליניארי. \\
\hline
הצמדת מטען $C$ ושיקוף פריטי $P$ & לרוב אוניטריים &
אך בצירופים מסוימים עם $T$ נוצרת אנטי־אוניטריות. \\
\hline
\end{tabular}
\end{center}

\subsection*{דוגמה: היפוך זמן עבור מצב תנע}

אם $T$ הוא אופרטור היפוך זמן, אז:
\[
T\, i \, T^{-1} = -i, \qquad
T\, \hat{p} \, T^{-1} = -\hat{p}, \qquad
T\, \hat{x} \, T^{-1} = \hat{x}.
\]
כדי ש-$i$ ישתנה ל-$-i$, על $T$ להיות אנטי־ליניארי.

\section*{סיכום עיקרי}

\begin{center}
\begin{tabular}{|c|c|c|c|p{5cm}|}
\hline
\textbf{מושג} & \textbf{ליניאריות} & \textbf{שימור סקלרים} & \textbf{שימור מכפלה פנימית} & \textbf{שימוש פיזיקלי} \\
\hline
אוניטרי $U$ & ליניארי & $a \to a$ &
$\langle U\psi, U\phi \rangle = \langle \psi, \phi \rangle$ &
סיבובים, הזזות, סימטריות מרחביות. \\
\hline
אנטי־אוניטרי $A$ & אנטי־ליניארי & $a \to a^*$ &
$\langle A\psi, A\phi \rangle = \langle \phi, \psi \rangle$ &
היפוך זמן, הפרות $CP$, וכו׳. \\
\hline
\end{tabular}
\end{center}

\subsection*{\LR{Takeaways}}

\begin{enumerate}
    \item אנטי־ליניאריות פירושה צימוד של הסקלרים — מנגנון מתמטי הכרחי לייצוג היפוך כיווניות בזמן.
    \item אנטי־אוניטריות משמרת נורמות אך הופכת את כיווניות הפאזה.
    \item משפט ויגנר הוא הגשר בין סימטריות פיזיקליות לאופרטורים אוניטריים או אנטי־אוניטריים.
    \item כל סימטריה קוונטית אמיתית מיוצגת באחד משני האופנים הללו — אין אחרים עקב הדרישה לשימור הסתברויות.
\end{enumerate}
